\section{Numerically Solving the 2-Dimensional System}

Now, we apply our methods in the previous section to solve the stochastic
FitzHugh--Nagumo system.  We use \(u\) for the activator field and \(v\) for
the inhibitor field throughout:
\begin{align}
  \frac{\partial u}{\partial t}
    &= D_u \nabla^2 u
      + \frac{1}{\tau_u}(1-u^2)(u-v)
      + \eta(\mathbf r,t),
      \label{eq:fhn-u}\\
  \frac{\partial v}{\partial t}
    &= D_v \nabla^2 v
      + \frac{1}{\tau_v}(u-\alpha v+\beta).
      \label{eq:fhn-v}
\end{align}
Here, \(D_u\) and \(D_v\) are the diffusion constants, while \(\tau_u\)
and \(\tau_v\) are the characteristic response times of the activator and
inhibitor, respectively.  The parameters \(\alpha\) and \(\beta\) determine
the relaxation and homogeneous steady-state value of the inhibitor.  The
noise field \(\eta\) acts only on the activator.  The one-dimensional
simulations considered in the accompanying work are obtained by replacing
\(\nabla^2\) with \(\partial_\ell^2\) and imposing periodic boundary
conditions.

\subsection{Attempting to Linearize the Non-linearity}

In all the algebra that will follow, it is easy to lose track of our goal.
As such, we will explicitly state it as finding an update rule that
satisfies
\begin{equation}
  \boldsymbol{\psi}^{n+1}
    = \mathcal M\!\left(\boldsymbol{\psi}^n,
      \boldsymbol{\eta}^{n+1}\right),
  \qquad
  \boldsymbol{\psi}^n \equiv
  \begin{bmatrix}
    \boldsymbol u^n\\
    \boldsymbol v^n
  \end{bmatrix},
  \label{eq:update-goal}
\end{equation}
where the upper index \(n\) characterizes the current time step and bold
symbols denote vectors containing the values at all spatial grid points.
Because the activator equation is non-linear and stochastic,
\(\mathcal M\) is not a single constant matrix.  Nevertheless, each
deterministic Crank--Nicolson step can be reduced to a linear solve and a
fixed-point iteration.

We see that the difference between Eq.~\eqref{eq:fhn-u} and the simple
diffusion equation is the non-linearity.  For reasons that will later
become clear, we let it be represented by the function
\begin{equation}
  \varphi(t,u,v)
    \equiv \frac{1}{\tau_u}(1-u^2)(u-v).
  \label{eq:phi-definition}
\end{equation}
Despite the problems that this presents, it is worth noting now that
\(\varphi\) is easy to compute given full knowledge of \(u\) and \(v\) at
a given time.  We first omit the noise term and rewrite the deterministic
part of the equations as
\begin{align}
  \frac{\partial u}{\partial t}
    &= D_u\nabla^2u+\varphi,
    \label{eq:deterministic-u}\\
  \frac{\partial v}{\partial t}
    &= D_v\nabla^2v+\frac{1}{\tau_v}(u-\alpha v+\beta).
    \label{eq:deterministic-v}
\end{align}
The noise will be applied after this deterministic step, in the same
operator-split form used by the simulations.

Let \(\mathbf L\) be the matrix that approximates the Laplacian on the
spatial grid.  For the current one-dimensional periodic simulations,
\begin{equation}
  \mathbf L = \frac{1}{a^2}(\mathbf G-2\mathbf I),
  \label{eq:periodic-laplacian}
\end{equation}
where \(a\) is the lattice spacing and \(\mathbf G\) has entries equal to
one for the two neighboring grid points, including the two corner entries
that enforce the periodic boundary condition.  For the triangulated
spherical surface, \(\mathbf L\) is replaced by the graph-based Laplacian
described in the Appendix.  The derivation below is unchanged once
\(\mathbf L\) has been chosen.

Applying Crank--Nicolson to the deterministic equations gives
\begin{align}
  \frac{\overline{\boldsymbol u}^{\,n+1}-\boldsymbol u^n}{\Delta t}
    &=
      \frac{D_u}{2}\mathbf L
      \left(\overline{\boldsymbol u}^{\,n+1}+\boldsymbol u^n\right)
      +\frac{1}{2}
      \left(\boldsymbol\varphi^{n+1}
      +\boldsymbol\varphi^n\right),
      \label{eq:cn-u}\\
  \frac{\overline{\boldsymbol v}^{\,n+1}-\boldsymbol v^n}{\Delta t}
    &=
      \frac{D_v}{2}\mathbf L
      \left(\overline{\boldsymbol v}^{\,n+1}+\boldsymbol v^n\right)
      +\frac{1}{2\tau_v}
      \left[
        \overline{\boldsymbol u}^{\,n+1}+\boldsymbol u^n
        -\alpha\left(
          \overline{\boldsymbol v}^{\,n+1}+\boldsymbol v^n
        \right)
        +2\beta\boldsymbol 1
      \right].
      \label{eq:cn-v}
\end{align}
The overline denotes the result of the deterministic substep, before the
noise is added.  We define
\begin{align}
  \mathbf A_u
    &\equiv \mathbf I-\frac{D_u\Delta t}{2}\mathbf L,
  &\mathbf B_u
    &\equiv \mathbf I+\frac{D_u\Delta t}{2}\mathbf L,
    \label{eq:uv-matrices-u}\\
  \mathbf A_v
    &\equiv \mathbf I-\frac{D_v\Delta t}{2}\mathbf L
      +\frac{\alpha\Delta t}{2\tau_v}\mathbf I,
  &\mathbf B_v
    &\equiv \mathbf I+\frac{D_v\Delta t}{2}\mathbf L
      -\frac{\alpha\Delta t}{2\tau_v}\mathbf I,
    \label{eq:uv-matrices-v}\\
  \mathbf C
    &\equiv \frac{\Delta t}{2}\mathbf I.
  &&
  \label{eq:c-matrix}
\end{align}
With these definitions, Eqs.~\eqref{eq:cn-u} and \eqref{eq:cn-v} can be
put into matrix form:
\begin{equation}
  \begin{bmatrix}
    \mathbf A_u & \mathbf 0\\
    -\mathbf C/\tau_v & \mathbf A_v
  \end{bmatrix}
  \begin{bmatrix}
    \overline{\boldsymbol u}^{\,n+1}\\
    \overline{\boldsymbol v}^{\,n+1}
  \end{bmatrix}
  =
  \begin{bmatrix}
    \mathbf B_u & \mathbf 0\\
    \mathbf C/\tau_v & \mathbf B_v
  \end{bmatrix}
  \begin{bmatrix}
    \boldsymbol u^n\\
    \boldsymbol v^n
  \end{bmatrix}
  +
  \begin{bmatrix}
    \mathbf C\left(
      \boldsymbol\varphi^n+\boldsymbol\varphi^{n+1}
    \right)\\
    (\beta\Delta t/\tau_v)\boldsymbol 1
  \end{bmatrix}.
  \label{eq:block-system}
\end{equation}
We see that we could not put the problem into the constant-matrix form
that we originally wanted, but as it turns out, this is good enough.

Assuming that \(\boldsymbol\varphi^{n+1}\) is a fixed point, we can use
fixed-point iteration to find a stable solution.  Let
\(\boldsymbol\varphi^{n+1,0}=\boldsymbol\varphi^n\) be our initial guess.
At iteration \(k\), we solve
\begin{equation}
  \begin{bmatrix}
    \mathbf A_u & \mathbf 0\\
    -\mathbf C/\tau_v & \mathbf A_v
  \end{bmatrix}
  \begin{bmatrix}
    \overline{\boldsymbol u}^{\,n+1,k+1}\\
    \overline{\boldsymbol v}^{\,n+1,k+1}
  \end{bmatrix}
  =
  \begin{bmatrix}
    \mathbf B_u & \mathbf 0\\
    \mathbf C/\tau_v & \mathbf B_v
  \end{bmatrix}
  \begin{bmatrix}
    \boldsymbol u^n\\
    \boldsymbol v^n
  \end{bmatrix}
  +
  \begin{bmatrix}
    \mathbf C\left(
      \boldsymbol\varphi^n+\boldsymbol\varphi^{n+1,k}
    \right)\\
    (\beta\Delta t/\tau_v)\boldsymbol 1
  \end{bmatrix}.
  \label{eq:fixed-point-system}
\end{equation}
We then use \(\overline{\boldsymbol u}^{\,n+1,k+1}\) and
\(\overline{\boldsymbol v}^{\,n+1,k+1}\) to calculate
\(\boldsymbol\varphi^{n+1,k+1}\) from Eq.~\eqref{eq:phi-definition}.  We
repeat until successive values of \(\boldsymbol\varphi\) converge.  In the
current simulations, the iteration stops when their norm differs by less
than \(10^{-10}\), with a maximum of \(100\) iterations.

\subsection{Active Noise in the Simulations}

The active noise in Eq.~\eqref{eq:fhn-u} is an
Ornstein--Uhlenbeck process with zero mean and autocovariance
\begin{equation}
  \left\langle
    \eta(\ell,t)\eta(\ell',t')
  \right\rangle
  =
  \delta(\ell-\ell')\frac{A}{\tau_c}
  \exp\left(-\frac{|t-t'|}{\tau_c}\right).
  \label{eq:noise-covariance}
\end{equation}
Here, \(A\) is the integrated noise strength and \(\tau_c\) is the
correlation time.  Equivalently, the noise field satisfies
\begin{equation}
  \frac{\partial\eta}{\partial t}
    =-\frac{\eta}{\tau_c}
      +\frac{\sqrt{2A}}{\tau_c}\xi(\ell,t),
  \label{eq:ou-equation}
\end{equation}
where \(\xi\) is a Gaussian white-noise field.  We discretize this equation
using Euler--Maruyama:
\begin{equation}
  \eta_i^{n+1}
    =\left(1-\frac{\Delta t}{\tau_c}\right)\eta_i^n
      +\frac{\sqrt{2A\Delta t}}{\tau_c}\xi_i^n,
  \label{eq:ou-update}
\end{equation}
where each \(\xi_i^n\) is an independent standard Gaussian random
variable.  After the deterministic fixed-point iteration has converged,
the operator-split noise update is
\begin{equation}
  \boldsymbol u^{n+1}
    =\overline{\boldsymbol u}^{\,n+1}
      +\Delta t\,\boldsymbol\eta^{n+1},
  \qquad
  \boldsymbol v^{n+1}
    =\overline{\boldsymbol v}^{\,n+1}.
  \label{eq:split-noise-update}
\end{equation}
When \(\tau_c=0\), the simulations use the Gaussian white-noise limit
directly:
\begin{equation}
  \boldsymbol u^{n+1}
    =\overline{\boldsymbol u}^{\,n+1}
      +\sqrt{2A\Delta t}\,\boldsymbol\xi^n.
  \label{eq:white-noise-update}
\end{equation}

For comparison with the numerical results, time is measured in units for
which \(\tau_u=1\).  The simulations use
\(D_u=1\), \(D_v=5\), \(\alpha=0.5\), and \(\beta=0.7\), while
\(\tau_v\), \(\tau_c\), and \(A\) are varied.  The homogeneous quiescent
initial state is
\begin{equation}
  u_0=-1,
  \qquad
  v_0=\frac{\beta-1}{\alpha}=-0.6.
  \label{eq:quiescent-state}
\end{equation}
For the current periodic grid, the lattice spacing is \(a=0.1\), and
simulation snapshots are recorded every \(0.1\) time units.
